\chapter{Sinusoids, Phasors, and Complex Exponentials}

Chapter~3 showed that frequency analysis is projecting a signal onto an orthogonal basis
of sinusoids. This chapter examines those sinusoids in detail.
Almost everything we do in this book -- the Fourier series, the Fourier transform, filters, the
frequency response of a system, sampling, modulation -- is built on a single wave that goes up
and down steadily at one pure rate. In this chapter we look at that wave from three angles at
once: as a wiggle in time, as the shadow of uniform circular motion, and as a single rotating
complex number that carries amplitude, frequency, and phase together. That rotating-vector view
simplifies nearly every calculation in the field. By the end of the chapter, $e^{\jj \omega t}$
should feel like the most natural way to write down steady rotation.

\section{What a sinusoid is, and the three numbers that pin it down}
\label{ch04:three-numbers}

Start with the most familiar curve in mathematics, the cosine. In its bare form $\cos(t)$ it
oscillates between $+1$ and $-1$, completing one full cycle every $2\pi$ units of $t$. A real
signal in the world -- a sound pressure, a voltage on a wire, the side-to-side sway of your
forearm -- is rarely so plain. To describe an actual oscillation we need to stretch it, speed it
up, and slide it sideways. Those three operations correspond to exactly three numbers.

\begin{definition}[Continuous-time sinusoid]
\label{ch04:def-sinusoid}
A (real) continuous-time sinusoid is any signal of the form
\[
  x(t) \;=\; A\cos(\omega t + \phi),
\]
where $A \ge 0$ is the \emph{amplitude}, $\omega$ is the \emph{angular frequency} in radians per
second, and $\phi$ is the \emph{phase} in radians. The ordinary frequency is $f = \omega/(2\pi)$
in hertz (cycles per second), and the \emph{period} is $T = 1/f = 2\pi/\omega$ seconds.
\end{definition}

Let us take the three numbers one at a time, because each one has a clean geometric meaning that
will pay off later.

\textbf{Amplitude $A$} is a vertical scaling. It is how far the wave reaches above and below its
center line; the peak-to-peak swing is $2A$. Double $A$ and you double the loudness of a tone,
the brightness of a flicker, the size of a sway. Amplitude has the same physical units as the
signal itself -- volts, pascals, $g$ of acceleration.

\textbf{Angular frequency $\omega$} is a horizontal scaling: it controls how fast the argument of
the cosine runs. Because cosine repeats every time its argument advances by $2\pi$, the signal
completes a full cycle whenever $\omega t$ grows by $2\pi$, i.e.\ every $T = 2\pi/\omega$ seconds.
Bigger $\omega$ means the cycles are squeezed closer together in time -- a higher pitch, a faster
wiggle.

\textbf{Phase $\phi$} is a horizontal shift. Writing $A\cos(\omega t + \phi) =
A\cos\!\big(\omega(t + \phi/\omega)\big)$ shows that adding $\phi$ slides the whole wave earlier
in time by $\phi/\omega$ seconds. Phase tells you \emph{where in its cycle} the wave is at the
instant $t = 0$. A cosine has $\phi = 0$ (it starts at its peak); a sine is a cosine with
$\phi = -\pi/2$, because $\sin(\omega t) = \cos(\omega t - \pi/2)$.

\begin{figure}[H]
\centering
\begin{tikzpicture}
\begin{axis}[
    width=11.5cm, height=5.5cm,
    title={The three parameters of a sinusoid: amplitude, frequency, and phase},
    title style={font=\small},
    xlabel={time $t$ (s)}, ylabel={$x(t)$},
    xmin=0, xmax=2, ymin=-2.4, ymax=2.4,
    axis lines=left, samples=180, domain=0:2,
    xlabel style={at={(axis description cs:0.5,-0.10)}, anchor=north},
    legend style={font=\tiny, at={(axis description cs:0.5,-0.28)}, anchor=north, legend columns=2},
    every axis plot/.append style={very thick},
    grid=both, grid style={gray!18},
]
\addplot[sigblue] {1.5*cos(deg(2*pi*1*x))};
\addlegendentry{$1.5\cos(2\pi t)$}
\addplot[specred,dashed] {1.0*cos(deg(2*pi*1*x))};
\addlegendentry{$1.0\cos(2\pi t)$ (smaller $A$)}
\addplot[quotientorange,densely dotted] {1.5*cos(deg(2*pi*2*x))};
\addlegendentry{$1.5\cos(2\pi\cdot 2 t)$ (higher $\omega$)}
\addplot[structurepurple,dash dot] {1.5*cos(deg(2*pi*1*x + 90))};
\addlegendentry{$1.5\cos(2\pi t + \pi/2)$ (phase)}
\end{axis}
\end{tikzpicture}
\caption{The three knobs. Amplitude scales the height, angular frequency squeezes the cycles
together, and phase slides the wave sideways. Each curve changes exactly one parameter from the
solid blue reference.}
\label{ch04:fig-three-knobs}
\end{figure}

\begin{trapbox}
Do not confuse $f$ (hertz, cycles per second) with $\omega$ (radians per second). They differ by
the factor $2\pi$ that converts one full turn into one cycle: $\omega = 2\pi f$. Mixing them up is
the single most common unit error in DSP -- a filter that should cut at $50$ Hz ends up cutting at
$50/(2\pi) \approx 8$ Hz, or at $50\cdot 2\pi \approx 314$ Hz. Always write the units, and when in
doubt ask: ``does one cycle make my argument advance by $2\pi$?''
\end{trapbox}

\begin{intuitionbox}
A sinusoid is the \emph{only} waveform that keeps its shape under every operation we care about.
Differentiate it, integrate it, delay it, run it through any linear time-invariant system, and out
comes another sinusoid of the \emph{same frequency} -- only $A$ and $\phi$ change. Square waves,
triangle waves, your own EMG envelope: none of them have this property. This fact is why
the theory of frequency is built out of sinusoids. We will make it precise
in Section~\ref{ch04:eigen}.
\end{intuitionbox}

\section{Why cosine alone is awkward: the addition problem}
\label{ch04:addition}

A small task motivates the rest of the chapter. Suppose two sources produce
sinusoids of the \emph{same} frequency $\omega$ but different amplitudes and phases, and they add
together (two speakers, two reflections of a radio wave, two muscle fibers firing). What is the
sum?
\[
  A_1\cos(\omega t + \phi_1) + A_2\cos(\omega t + \phi_2) \;=\; ?
\]
We \emph{know} the answer must again be a sinusoid of frequency $\omega$, because adding two
same-frequency cosines cannot manufacture a new rate of oscillation. So the result is
$A_3\cos(\omega t + \phi_3)$ for some new amplitude and phase. But finding $A_3$ and $\phi_3$ with
trig identities is genuinely painful. Expand each term with the angle-addition formula
$\cos(\omega t + \phi) = \cos\phi\cos\omega t - \sin\phi\sin\omega t$, collect the
$\cos\omega t$ and $\sin\omega t$ pieces, and then re-bundle them back into a single cosine. It
works, but it is bookkeeping with no insight, and it gets worse with three or four terms.

\begin{walkbox}
Before reading on, try the special case $\cos(\omega t) + \cos(\omega t + \pi/2)$ by hand with
trig. You should get $\sqrt{2}\,\cos(\omega t + \pi/4)$. Notice how much algebra that took for two
unit-amplitude terms. Now imagine summing the hundred sinusoids that make up a violin note. We
need a better tool.
\end{walkbox}

The better tool is to treat a sinusoid as the shadow of a rotating point rather than as an
isolated wiggle.

\section{Rotation on the unit circle, and where cosine comes from}
\label{ch04:rotation}

Picture a point moving counterclockwise around a circle of radius $A$ at a steady angular speed
$\omega$. At time $t$ its angle (measured from the positive horizontal axis) is $\omega t + \phi$,
where $\phi$ is wherever it started at $t = 0$. Its coordinates are, by the definition of sine and
cosine,
\[
  \big(\,A\cos(\omega t + \phi)\,,\; A\sin(\omega t + \phi)\,\big).
\]
Consider the first coordinate. \emph{The horizontal position of a steadily rotating point is
exactly our sinusoid.} A cosine can therefore be viewed as the
\emph{shadow} that a rotating point casts on the horizontal axis. The vertical shadow is a sine.
The wiggle in time that we plotted in Figure~\ref{ch04:fig-three-knobs} is what you get when you
unroll that circular motion against a moving time axis.

\begin{figure}[H]
\centering
\begin{tikzpicture}[scale=1.0]
% --- unit circle on the left ---
\begin{scope}[shift={(0,0)}]
  \draw[gray!35,->] (-1.7,0) -- (1.9,0) node[right]{$\mathrm{Re}$};
  \draw[gray!35,->] (0,-1.7) -- (0,1.9) node[above]{$\mathrm{Im}$};
  \draw[sigblue,thick] (0,0) circle (1.3);
  % rotating vector at 50 degrees
  \draw[->,very thick,structurepurple] (0,0) -- (50:1.3) node[midway,above left]{$A$};
  \fill[structurepurple] (50:1.3) circle (1.3pt);
  % angle arc
  \draw[quotientorange,->] (0.55,0) arc (0:50:0.55);
  \node[quotientorange] at (28:0.78) {\small $\omega t{+}\phi$};
  % projection lines
  \draw[specred,densely dashed] (50:1.3) -- ({1.3*cos(50)},0);
  \fill[specred] ({1.3*cos(50)},0) circle (1.3pt);
  \draw[targetgreen,densely dashed] (50:1.3) -- (0,{1.3*sin(50)});
  \fill[targetgreen] (0,{1.3*sin(50)}) circle (1.3pt);
  \node[specred,below] at ({1.3*cos(50)},-0.05) {\small $A\cos$};
  \node[targetgreen,left] at (-0.05,{1.3*sin(50)}) {\small $A\sin$};
\end{scope}
% --- time plot on the right ---
\begin{scope}[shift={(4.6,0)}]
  \draw[gray!35,->] (0,-1.6) -- (0,1.7) node[above]{$x$};
  \draw[gray!35,->] (-0.2,0) -- (4.6,0) node[right]{$t$};
  \draw[specred,thick,domain=0:4.4,samples=120,smooth,variable=\u]
        plot ({\u}, {1.3*cos(deg(\u*1.8)+50)});
  % link the starting height
  \draw[specred,densely dotted] (-3.5,{1.3*cos(50)}) -- (0,{1.3*cos(50)});
  \fill[specred] (0,{1.3*cos(50)}) circle (1.3pt);
\end{scope}
\end{tikzpicture}
\caption{Left: a vector of length $A$ rotates counterclockwise; at time $t$ its angle is
$\omega t + \phi$. Its projection onto the real axis (red) is $A\cos(\omega t+\phi)$; onto the
imaginary axis (green) is $A\sin(\omega t+\phi)$. Right: unrolling the red projection against time
gives the cosine wave.}
\label{ch04:fig-circle-to-wave}
\end{figure}

\begin{happeningbox}
When you watch a cosine oscillate, what is \emph{really} happening is that something is going
around in a circle at constant speed and you are only allowed to see one coordinate of it. The
``slowing down'' near the peaks and the ``speeding through'' near zero -- the non-uniform pace of the
cosine curve -- is an effect of the projection. The
rotation itself is perfectly uniform. All the apparent acceleration lives in the geometry of the
shadow, not in the motion.
\end{happeningbox}

This reframing already explains the period. The rotating point returns to its start after one full
turn, $2\pi$ radians, which at speed $\omega$ takes $T = 2\pi/\omega$ seconds. It explains
amplitude: the radius of the circle. And it explains phase: the starting angle. All three numbers
describe one circle. The only thing missing is a way to write ``a rotating vector'' as a
single algebraic object so we can compute with it. For that we need the complex plane.

\section{Euler's formula: rotation written as one number}
\label{ch04:euler}

A point in the plane with coordinates $(a, b)$ can be packaged as a single complex number
$z = a + \jj b$, where $\jj$ is the imaginary unit with $\jj^2 = -1$. (Engineers write $\jj$ rather
than $i$, since $i$ is reserved for electric current.) Addition of complex numbers is just
vector addition of the points. The important operation is \emph{multiplication}: multiplying by a complex
number of magnitude $1$ and angle $\theta$ rotates the plane by $\theta$. So complex multiplication
\emph{is} rotation -- which is exactly the operation we needed for sinusoids.

The bridge between the additive picture (coordinates) and the multiplicative picture (rotation) is
Euler's formula, the most important single equation in this book.

\begin{theorem}[Euler's formula]
\label{ch04:thm-euler}
For any real angle $\theta$,
\[
  e^{\jj\theta} \;=\; \cos\theta + \jj\sin\theta .
\]
Consequently $e^{\jj\theta}$ is the point on the unit circle at angle $\theta$, and
$\abs{e^{\jj\theta}} = \sqrt{\cos^2\theta + \sin^2\theta} = 1$.
\end{theorem}

\begin{whybox}
Why should the exponential function, born from compound interest and growth, have anything to do
with circles? Because both are defined by the same differential law, just with a twist of $\jj$.
The real exponential satisfies $\frac{\dd}{\dd t}e^{t} = e^{t}$: its rate of change points the
\emph{same} direction as its value, so it grows. Now consider $u(t) = e^{\jj t}$. Differentiating,
$u'(t) = \jj\, e^{\jj t} = \jj\, u(t)$. Multiplying by $\jj$ is a $90^\circ$ rotation, so the
velocity of the point $u(t)$ is always exactly perpendicular to its position vector. A velocity
forever perpendicular to the radius, with constant speed, is precisely the definition of uniform
circular motion. The exponential whose rate is rotated by a quarter turn cannot grow outward; it
can only \emph{circle}. That is Euler's formula, read as a sentence about motion.
\end{whybox}

\begin{focusbox}
A second, algebraic, route. Write down the Taylor series of $e^{x}$ and substitute $x = \jj\theta$:
\[
  e^{\jj\theta} = \sum_{n=0}^{\infty}\frac{(\jj\theta)^n}{n!}
   = \underbrace{\Big(1 - \tfrac{\theta^2}{2!} + \tfrac{\theta^4}{4!} - \cdots\Big)}_{\cos\theta}
   + \jj\underbrace{\Big(\theta - \tfrac{\theta^3}{3!} + \tfrac{\theta^5}{5!} - \cdots\Big)}_{\sin\theta},
\]
using $\jj^2 = -1$, $\jj^3 = -\jj$, $\jj^4 = 1$ to sort the powers into real and imaginary buckets.
The even powers reassemble the cosine series; the odd powers, each carrying one factor of $\jj$,
reassemble the sine series. The companion calculus book derives these three series; here we only
note how they fit together.
\end{focusbox}

Now feed time into the angle. A point rotating at angular speed $\omega$, starting at angle $\phi$
on a circle of radius $A$, is the single complex-valued signal
\[
  z(t) \;=\; A\,e^{\jj(\omega t + \phi)} \;=\; \underbrace{A e^{\jj\phi}}_{\text{constant}}\; e^{\jj\omega t}.
\]
This is the \textbf{complex exponential} (also called a complex sinusoid). It contains all three
parameters. Its real part is our cosine, its imaginary part is our sine:
\[
  \mathrm{Re}\,z(t) = A\cos(\omega t + \phi), \qquad \mathrm{Im}\,z(t) = A\sin(\omega t + \phi).
\]
And it factors neatly: the time-dependence $e^{\jj\omega t}$ (the spinning) separates from a
constant $Ae^{\jj\phi}$ that carries the amplitude and the starting angle. That constant has a name.

Figure~\ref{ch04:fig-helix} visualizes this geometry in three dimensions: as time advances, the
spinning vector sweeps out a 3D helix whose shadows on the coordinate walls are the real cosine and
imaginary sine waves.

\begin{figure}[H]
\centering
\begin{tikzpicture}[
  x={(0.95cm, -0.15cm)},
  y={(0.35cm, 0.45cm)},
  z={(0cm, 1.05cm)},
  scale=1.05, >=Stealth
]
  % Axes
  \draw[->, gray!60, thick] (0,0,0) -- (7.5,0,0) node[below right, font=\small] {time $t$};
  \draw[->, gray!60, thick] (0,-1.4,0) -- (0,1.8,0) node[above left, font=\small] {$\operatorname{Re}$};
  \draw[->, gray!60, thick] (0,0,-1.4) -- (0,0,1.6) node[above, font=\small] {$\operatorname{Im}$};

  % Unit circle at t = 0
  \draw[gray!40, dashed] plot[domain=0:360, samples=60] (0, {cos(\x)}, {sin(\x)});

  % Projections onto coordinate planes
  % 1. Horizontal plane (t, Re): real part = cos(omega t)
  \draw[specred, thick, domain=0:6.28, samples=100, smooth]
    plot (\x, {cos(deg(2*\x))}, 0);
  \node[specred, font=\footnotesize, right] at (6.3, 1.0, 0) {$\operatorname{Re}\{e^{\jj\omega t}\} = \cos(\omega t)$};

  % 2. Vertical plane (t, Im): imaginary part = sin(omega t)
  \draw[targetgreen, thick, domain=0:6.28, samples=100, smooth]
    plot (\x, 0, {sin(deg(2*\x))});
  \node[targetgreen, font=\footnotesize, above right] at (6.3, 0, 0.9) {$\operatorname{Im}\{e^{\jj\omega t}\} = \sin(\omega t)$};

  % 3. The 3D Helix: z(t) = e^{j omega t} = (t, cos, sin)
  \draw[structurepurple, line width=1.6pt, domain=0:6.28, samples=140, smooth]
    plot (\x, {cos(deg(2*\x))}, {sin(deg(2*\x))});
  \node[structurepurple, font=\small\bfseries, above] at (3.14, 0, 1.35) {3D Helix: $e^{\jj\omega t}$};

  % Snapshot at t0 = 1.0
  \coordinate (O0) at (1.0, 0, 0);
  \coordinate (P) at (1.0, {cos(deg(2*1.0))}, {sin(deg(2*1.0))});
  \coordinate (P_re) at (1.0, {cos(deg(2*1.0))}, 0);
  \coordinate (P_im) at (1.0, 0, {sin(deg(2*1.0))});

  \draw[->, sigblue, very thick] (O0) -- (P) node[pos=0.6, right=3pt, font=\scriptsize, fill=white, inner sep=1pt] {phasor};
  \fill[sigblue] (P) circle (1.8pt);
  \draw[gray!60, densely dashed] (P) -- (P_re);
  \draw[gray!60, densely dashed] (P) -- (P_im);
  \fill[specred] (P_re) circle (1.4pt);
  \fill[targetgreen] (P_im) circle (1.4pt);

  % Cylindrical guide circle at t = 2*pi
  \draw[gray!30, dotted] plot[domain=0:360, samples=50] (6.28, {cos(\x)}, {sin(\x)});
\end{tikzpicture}
\caption{The complex exponential helix $e^{\jj\omega t}$ and its orthogonal projections. As time $t$ increases, the rotating unit phasor traces a 3D corkscrew in $(t, \operatorname{Re}, \operatorname{Im})$ space. Casting a shadow onto the horizontal real plane $(t, \operatorname{Re})$ produces the cosine wave $\cos(\omega t)$ (red). Casting a shadow onto the vertical imaginary plane $(t, \operatorname{Im})$ produces the sine wave $\sin(\omega t)$ (green). Looking straight down the time axis collapses the entire 3D helix onto the 2D unit circle.}
\label{ch04:fig-helix}
\end{figure}

\section{Phasors: freezing the spin}
\label{ch04:phasors}

When every sinusoid in a problem shares the same frequency $\omega$, the factor $e^{\jj\omega t}$
is common to all of them -- they all spin together, in lockstep. So we can mentally ride along with
the rotation, ``stop the clock,'' and keep only the part that distinguishes one sinusoid from
another: its amplitude and phase. That frozen complex number is the phasor.

\begin{definition}[Phasor]
\label{ch04:def-phasor}
The \emph{phasor} of the sinusoid $A\cos(\omega t + \phi)$ is the complex number
\[
  X \;=\; A\,e^{\jj\phi} \;=\; A\cos\phi + \jj A\sin\phi .
\]
The sinusoid is recovered as $x(t) = \mathrm{Re}\big(X\,e^{\jj\omega t}\big)$. The magnitude
$\abs{X} = A$ is the amplitude; the angle $\arg X = \phi$ is the phase.
\end{definition}

A phasor is an arrow in the complex plane. Its length is the amplitude, its direction is the phase.
To get the actual signal back, you set this arrow spinning at $\omega$ and watch its real-axis
shadow. Same-frequency sinusoids become \emph{static} arrows, and
\textbf{adding sinusoids becomes adding arrows.}

\begin{figure}[H]
\centering
\begin{tikzpicture}[scale=1.25,>=Stealth]
  \draw[gray!35,->] (-0.6,0) -- (3.0,0) node[right]{$\mathrm{Re}$};
  \draw[gray!35,->] (0,-0.6) -- (0,2.3) node[above]{$\mathrm{Im}$};
  % phasor 1: A1=2, phi1=0
  \draw[->,very thick,sigblue] (0,0) -- (2,0);
  \node[sigblue,below] at (1.0,-0.05) {\small $X_1 = 2$};
  % phasor 2: A2=1.4, phi2=90 -> (0,1.4), drawn tip-to-tail from end of X1
  \draw[->,very thick,specred] (2,0) -- (2,1.4);
  \node[specred,right] at (2.02,0.7) {\small $X_2 = 1.4\,e^{\jj\pi/2}$};
  % resultant
  \draw[->,very thick,structurepurple] (0,0) -- (2,1.4);
  \node[structurepurple,above left] at (1.0,0.9) {\small $X_3 = X_1{+}X_2$};
  % angle of resultant
  \draw[quotientorange,->] (0.6,0) arc (0:35:0.6);
  \node[quotientorange] at (20:0.83) {\small $\phi_3$};
\end{tikzpicture}
\caption{Adding two same-frequency sinusoids is tip-to-tail vector addition of their phasors. Here
$2\cos\omega t + 1.4\cos(\omega t + \tfrac{\pi}{2})$ has phasor sum $X_3 = 2 + 1.4\jj$, so the
result is $\abs{X_3}\cos(\omega t + \phi_3)$ with $\abs{X_3} = \sqrt{2^2 + 1.4^2} \approx 2.44$ and
$\phi_3 = \arctan(1.4/2) \approx 0.61$ rad.}
\label{ch04:fig-phasor-add}
\end{figure}

This solves the awkward addition problem of Section~\ref{ch04:addition}. Because
$\mathrm{Re}$ is linear,
\[
  A_1\cos(\omega t+\phi_1) + A_2\cos(\omega t+\phi_2)
   = \mathrm{Re}\Big[(X_1 + X_2)\,e^{\jj\omega t}\Big],
\]
so we just add the phasors $X_1 + X_2 = X_3$ as complex numbers and read off the new amplitude and
phase as $\abs{X_3}$ and $\arg X_3$, with no trig identities.

\begin{example}[The walking-question sum, finally easy]
\label{ch04:ex-easy-sum}
Recompute $\cos(\omega t) + \cos(\omega t + \tfrac{\pi}{2})$. The phasors are $X_1 = 1\,e^{\jj 0} =
1$ and $X_2 = 1\,e^{\jj\pi/2} = \jj$. Their sum is $X_3 = 1 + \jj$, which has magnitude
$\abs{1+\jj} = \sqrt{2}$ and angle $\arg(1+\jj) = \pi/4$. Therefore
\[
  \cos(\omega t) + \cos\!\big(\omega t + \tfrac{\pi}{2}\big) = \sqrt{2}\,\cos\!\big(\omega t + \tfrac{\pi}{4}\big).
\]
This takes two lines, and the geometry shows \emph{why}: a horizontal unit arrow plus a vertical unit arrow makes a
diagonal arrow of length $\sqrt 2$ at $45^\circ$.
\end{example}

\begin{example}[Three-way sum]
\label{ch04:ex-three-sum}
Find $3\cos(\omega t) + 2\cos(\omega t - \tfrac{2\pi}{3}) + 2\cos(\omega t + \tfrac{2\pi}{3})$.
Phasors: $X_1 = 3$, $X_2 = 2e^{-\jj 2\pi/3} = 2(-\tfrac12 - \jj\tfrac{\sqrt3}{2}) = -1 - \jj\sqrt3$,
and $X_3 = 2e^{\jj 2\pi/3} = -1 + \jj\sqrt3$. Summing: $X_2 + X_3 = -2$ (the imaginary parts
cancel), so $X_1 + X_2 + X_3 = 3 - 2 = 1$. The result is simply $1\cdot\cos(\omega t)$. The two
side terms, symmetric about the real axis, partially cancel the main term. The same calculation with trig
identities takes much longer.
\end{example}

\begin{whybox}
Why are we allowed to add the phasors and ignore the spinning $e^{\jj\omega t}$? Because all the
terms share it. Factor it out: $\sum_k \mathrm{Re}(X_k e^{\jj\omega t}) =
\mathrm{Re}\big((\sum_k X_k)\,e^{\jj\omega t}\big)$. The common rotation is applied once, at the end,
to the summed arrow. This only works when the frequencies match -- if two sinusoids spin at
different rates, their arrows drift relative to each other and there is no single frozen picture.
Different frequencies do not add to a single sinusoid; that case leads to the Fourier series.
\end{whybox}

\begin{connectionbox}
Adding phasors is ordinary vector addition in $\R^2$. In the companion linear-algebra book, a
complex number $a + \jj b$ corresponds to the vector $(a,b)$, complex addition is componentwise
addition, and $\abs{\cdot}$ is the Euclidean norm. Complex arithmetic simply packages
two-dimensional vector geometry into compact algebraic notation. Later, when we project signals
onto sinusoids in the Fourier transform, the inner product $\inner{x}{y}$ from that book will be
the operation we apply -- and the complex exponentials will be the vectors we project onto.
\end{connectionbox}

\section{Two counter-rotating arrows: the negative frequency}
\label{ch04:negfreq}

We defined the phasor by taking the real part of a single rotating arrow. There is a second,
perfectly symmetric, way to recover a real cosine, and it will matter a great deal once we reach the
Fourier transform. Start from Euler's formula written both ways:
\[
  e^{\jj\theta} = \cos\theta + \jj\sin\theta, \qquad e^{-\jj\theta} = \cos\theta - \jj\sin\theta.
\]
Add them and the imaginary parts cancel; subtract them and the real parts cancel. Solving,
\[
  \boxed{\;\cos\theta = \frac{e^{\jj\theta} + e^{-\jj\theta}}{2}, \qquad
         \sin\theta = \frac{e^{\jj\theta} - e^{-\jj\theta}}{2\jj}.\;}
\]
Setting $\theta = \omega t$, a real cosine is the average of two arrows: one of length $\tfrac12$
spinning counterclockwise at $+\omega$, and one of length $\tfrac12$ spinning clockwise at
$-\omega$. At every instant their vertical (imaginary) components are mirror images and cancel,
leaving a purely real, horizontal sum that oscillates as $\cos\omega t$.

\begin{figure}[H]
\centering
\begin{tikzpicture}[scale=1.2,>=Stealth]
  \draw[gray!35,->] (-1.7,0) -- (1.9,0) node[right]{$\mathrm{Re}$};
  \draw[gray!35,->] (0,-1.7) -- (0,1.7) node[above]{$\mathrm{Im}$};
  \draw[gray!30] (0,0) circle (1.2);
  % +omega arrow at +40
  \draw[->,very thick,sigblue] (0,0) -- (40:1.2);
  \node[sigblue,anchor=south west] at (40:1.25) {\small $\tfrac12 e^{\jj\omega t}$};
  \draw[sigblue,->] (0.35,0) arc (0:40:0.35);
  % -omega arrow at -40
  \draw[->,very thick,specred] (0,0) -- (-40:1.2);
  \node[specred,anchor=north west] at (-40:1.25) {\small $\tfrac12 e^{-\jj\omega t}$};
  \draw[specred,->] (0.35,0) arc (0:-40:0.35);
  % vertical components canceling
  \draw[gray,densely dashed] (40:1.2) -- ({1.2*cos(40)},0);
  \draw[gray,densely dashed] (-40:1.2) -- ({1.2*cos(40)},0);
  % resultant on real axis
  \draw[->,very thick,structurepurple] (0,0) -- ({1.2*cos(40)},0);
  \node[structurepurple,anchor=south west] at (0.45,0.05) {\small $\cos\omega t$};
\end{tikzpicture}
\caption{A real cosine as two counter-rotating half-arrows. The blue one spins at $+\omega$, the
red at $-\omega$; their imaginary parts always cancel and their real parts add, tracing
$\cos\omega t$ on the horizontal axis. This is the origin of ``negative frequency.''}
\label{ch04:fig-counter-rotate}
\end{figure}

\begin{intuitionbox}
``Negative frequency'' does not mean less than zero in any physical sense. It is
simply \emph{the other direction of rotation}. Counterclockwise is positive, clockwise is negative.
A real-valued signal, having no imaginary part to keep track of a preferred sense of spin, must
always contain equal and opposite amounts of both -- which is why the spectrum of a real signal is
symmetric about zero. When we get to the Fourier transform you will see those mirrored peaks at
$+\omega$ and $-\omega$ and now you will know they are the two arrows of every real cosine.
\end{intuitionbox}

\begin{trapbox}
A single complex exponential $e^{\jj\omega t}$ has only \emph{one} frequency, $+\omega$, and its
spectrum is a single peak with nothing at $-\omega$. It is the real cosine, not the complex
exponential, that needs both. Quadrature receivers and the analytic signal of
Chapter~25 exploit exactly this: by carrying both a real and an imaginary
channel you can represent a one-sided spectrum and distinguish $+\omega$ from $-\omega$ -- something
a single real wire physically cannot do.
\end{trapbox}

\section{Complex exponentials are the eigenfunctions of LTI systems}
\label{ch04:eigen}

We promised in Section~\ref{ch04:three-numbers} that sinusoids are special because they survive
linear time-invariant (LTI) systems intact. Now we can state \emph{why}, and the reason is the
most important idea in the chapter: it explains why the Fourier transform is so useful.

A system is \emph{linear} if scaling and adding inputs scales and adds outputs, and
\emph{time-invariant} if delaying the input merely delays the output by the same amount (the
system's behavior does not depend on \emph{when} you poke it). Almost every device we model -- a
filter, a resonant cavity, an RC circuit, the mechanical response of soft tissue -- is LTI to good
approximation. Feed such a system the complex exponential $x(t) = e^{\jj\omega t}$. The claim is
that the output is the \emph{same} exponential, only scaled by a complex number that depends on
$\omega$:
\[
  e^{\jj\omega t} \;\xrightarrow{\;\text{LTI system } \mathcal{H}\;}\; H(\omega)\, e^{\jj\omega t}.
\]

\begin{whybox}
Here is the short proof. Let $y(t) = \mathcal{H}\{e^{\jj\omega t}\}$ be the output.
Time-invariance says that delaying the input by $\tau$ delays the output by $\tau$:
$\mathcal{H}\{e^{\jj\omega(t-\tau)}\} = y(t-\tau)$. But $e^{\jj\omega(t-\tau)} = e^{-\jj\omega\tau}\,
e^{\jj\omega t}$, and the scalar $e^{-\jj\omega\tau}$ pulls out by linearity, so the left side is
$e^{-\jj\omega\tau}\,y(t)$. Therefore $y(t-\tau) = e^{-\jj\omega\tau}\,y(t)$ for all $\tau$. A
function that, when shifted by any $\tau$, comes back as itself times $e^{-\jj\omega\tau}$ must be a
constant multiple of $e^{\jj\omega t}$. Writing that constant as $H(\omega)$ gives
$y(t) = H(\omega)e^{\jj\omega t}$. The exponential is preserved; only its complex amplitude is
multiplied.
\end{whybox}

In linear-algebra language: the complex exponential is an \emph{eigenfunction} of every LTI system,
and $H(\omega)$ is the corresponding \emph{eigenvalue}. The system does not bend or reshape the
exponential -- it can only stretch it (change amplitude, $\abs{H(\omega)}$) and rotate it (change
phase, $\arg H(\omega)$).

\begin{happeningbox}
This explains what audio engineers mean by a filter's ``frequency
response.'' The filter is a linear map, and exponentials are its eigenvectors. An eigenvector,
fed in, comes out merely scaled. So if you describe an arbitrary signal as a combination of
exponentials -- which is exactly what the Fourier transform does -- then the filter's effect is
trivial: multiply each exponential component by its own number $H(\omega)$. A hard problem
(``what does this circuit do to my complicated waveform?'') becomes a list of independent
multiplications. We decompose signals into sinusoids because sinusoids are the
coordinate system in which every LTI system is diagonal.
\end{happeningbox}

A single eigenvalue $H(\omega)$ is one complex number, but as $\omega$ ranges over all frequencies
it traces two real curves: a gain $\abs{H(\omega)}$ and a phase $\arg H(\omega)$. Together they are
the system's \emph{frequency response}. Figure~\ref{ch04:fig-freqresp} shows the pair for a simple
one-pole low-pass, $H(\omega) = 1/(1 + \jj\omega/\omega_c)$: low frequencies pass at near-unit gain
and near-zero phase, high frequencies are attenuated and delayed. Every LTI system you will meet is
described by one such gain curve and one such phase curve.

\begin{figure}[H]
\centering
\begin{tikzpicture}
\begin{axis}[
    width=6.6cm, height=4.6cm,
    title={Gain $\abs{H(\omega)}$},
    title style={font=\small},
    xlabel={$\omega/\omega_c$}, ylabel={$\abs{H(\omega)}$},
    xmin=0, xmax=5, ymin=0, ymax=1.1,
    axis lines=left, samples=160, domain=0:5,
    grid=both, grid style={gray!18},
]
\addplot[sigblue,very thick] {1/sqrt(1+x^2)};
\end{axis}
\begin{axis}[
    at={(7.6cm,0)},
    width=6.6cm, height=4.6cm,
    title={Phase $\arg H(\omega)$},
    title style={font=\small},
    xlabel={$\omega/\omega_c$}, ylabel={$\arg H(\omega)$ (rad)},
    xmin=0, xmax=5, ymin=-1.7, ymax=0.1,
    axis lines=left, samples=160, domain=0:5,
    grid=both, grid style={gray!18},
]
\addplot[specred,very thick] {-atan(x)*pi/180};
\end{axis}
\end{tikzpicture}
\caption{Frequency response of a one-pole low-pass $H(\omega)=1/(1+\jj\omega/\omega_c)$. Left: the
gain $\abs{H}$ falls from $1$ toward $0$, reaching $1/\sqrt2$ at $\omega=\omega_c$. Right: the phase
$\arg H = -\arctan(\omega/\omega_c)$ swings from $0$ down toward $-\pi/2$. Feeding in
$e^{\jj\omega t}$ returns the same exponential scaled by this one complex number.}
\label{ch04:fig-freqresp}
\end{figure}

For a real input cosine $\cos(\omega t) = \tfrac12 e^{\jj\omega t} + \tfrac12 e^{-\jj\omega t}$, the
same logic applies to the two counter-rotating half-arrows of Section~\ref{ch04:negfreq}. When the
system itself is real (a real input always produces a real output), conjugating
$\mathcal{H}\{e^{\jj\omega t}\} = H(\omega)e^{\jj\omega t}$ gives $H(-\omega) = \overline{H(\omega)}$.
So the $+\omega$ half-arrow is scaled by $H(\omega) = \abs{H(\omega)}e^{\jj\arg H(\omega)}$ and the
$-\omega$ half-arrow by its conjugate $\overline{H(\omega)} = \abs{H(\omega)}e^{-\jj\arg H(\omega)}$,
and their sum reassembles cleanly into a real cosine:
\[
  \cos(\omega t) \;\xrightarrow{\;\mathcal{H}\;}\;
  \tfrac12\abs{H(\omega)}e^{\jj(\omega t + \arg H(\omega))}
  + \tfrac12\abs{H(\omega)}e^{-\jj(\omega t + \arg H(\omega))}
  = \abs{H(\omega)}\,\cos\!\big(\omega t + \arg H(\omega)\big).
\]
Same frequency in, same frequency out -- amplitude scaled by the gain $\abs{H(\omega)}$, phase
shifted by $\arg H(\omega)$. This proves the claim made at the start of the chapter.

\begin{connectionbox}
``Eigenfunction,'' ``eigenvalue,'' ``diagonal in this basis'': these are the words of the
companion linear-algebra book, applied to vectors that happen to be functions. An LTI system is a
linear operator on the infinite-dimensional space of signals; the Fourier transform is the change
of basis that diagonalizes it. Chapters~12--14 build the convolution machinery that computes
$H(\omega)$ from a system's impulse response.
\end{connectionbox}

\section{Discrete-time complex exponentials (a first look)}
\label{ch04:discrete}

Real hardware samples. Your microcontroller does not see $\cos(\omega t)$ for continuous $t$; it
sees a sequence of numbers $x[n] = A\cos(\Omega n + \phi)$ taken every $T_s$ seconds, where $n$ is
an integer sample index and $\Omega = \omega T_s$ is the \emph{discrete} angular frequency in
radians \emph{per sample}. The complex version is $x[n] = A e^{\jj\phi} e^{\jj\Omega n}$ -- a point
that hops around the unit circle by $\Omega$ radians at each tick instead of sweeping continuously.

Two new effects appear as soon as we discretize, and both come straight from the geometry of hopping
around a circle. We mention them now and develop them fully in Chapters~5 and~18.

\textbf{Periodicity in frequency.} Because $e^{\jj 2\pi n} = 1$ for every integer $n$, increasing
$\Omega$ by $2\pi$ changes nothing: $e^{\jj(\Omega + 2\pi)n} = e^{\jj\Omega n}$. Two
discrete frequencies that differ by $2\pi$ produce \emph{identical} sample sequences. A continuous
sinusoid has a distinct identity at every frequency from $0$ to $\infty$; a discrete one repeats
every $2\pi$. All distinct discrete frequencies live in a single window of width $2\pi$, usually
taken as $(-\pi, \pi]$.

\textbf{Aliasing.} Two different continuous tones can land on the same samples. If $\Omega$ and
$\Omega + 2\pi$ give the same hops, then a fast continuous wave can masquerade as a slow one once
sampled. The highest distinct frequency is $\Omega = \pi$, half a turn per sample -- the Nyquist
limit, which Chapter~18 turns into the famous sampling theorem.

\begin{figure}[H]
\centering
\begin{tikzpicture}
\begin{axis}[
    width=12.0cm, height=5.2cm,
    title={A discrete-time sinusoid: sampling the continuous wave},
    title style={font=\small},
    xlabel={sample index $n$}, ylabel={$x[n]$},
    xlabel style={at={(axis description cs:0.5,-0.10)}, anchor=north},
    xmin=-0.4, xmax=12.4, ymin=-1.5, ymax=1.5,
    axis lines=left, grid=both, grid style={gray!18},
    legend style={font=\tiny, at={(axis description cs:0.5,-0.28)}, anchor=north, legend columns=2},
]
% underlying continuous wave Omega=0.5 rad/sample, plotted densely
\addplot[sigblue,thick,domain=0:12,samples=160] {cos(deg(0.5*x))};
\addlegendentry{$\cos(0.5\,t)$ underlying}
% samples as stems
\addplot[specred,ycomb,samples at={0,1,...,12}] {cos(deg(0.5*x))};
\addplot[specred,only marks,mark=*,mark size=1.6pt,samples at={0,1,...,12}] {cos(deg(0.5*x))};
\addlegendentry{samples $x[n]=\cos(0.5n)$}
\end{axis}
\end{tikzpicture}
\caption{A discrete-time sinusoid is samples of the rotating point taken at integer ticks. Here
$\Omega = 0.5$ rad/sample. Increasing $\Omega$ by $2\pi$ would leave every red dot exactly where it
is -- the source of aliasing and the periodicity of discrete frequency.}
\label{ch04:fig-discrete}
\end{figure}

Figure~\ref{ch04:fig-alias} makes the masquerade concrete. A slow tone and a much faster one,
chosen so their frequencies differ by exactly one full turn per sample, pass through the very same
integer samples. Looking only at the dots, you cannot tell which wave produced them.

\begin{figure}[H]
\centering
\begin{tikzpicture}
\begin{axis}[
    width=12.0cm, height=5.2cm,
    title={Aliasing: two continuous tones agreeing at all sample points},
    title style={font=\small},
    xlabel={time (samples)}, ylabel={value},
    xlabel style={at={(axis description cs:0.5,-0.14)}, anchor=north},
    xmin=-0.3, xmax=8.3, ymin=-1.5, ymax=1.5,
    axis lines=left, grid=both, grid style={gray!18},
    legend style={font=\tiny, at={(axis description cs:0.5,-0.34)}, anchor=north, legend columns=3},
]
% slow tone Omega = 0.4*pi rad/sample
\addplot[sigblue,thick,domain=0:8,samples=180] {cos(deg(0.4*pi*x))};
\addlegendentry{slow: $\Omega=0.4\pi$}
% fast tone Omega = 0.4*pi + 2*pi rad/sample, same samples
\addplot[quotientorange,densely dashed,domain=0:8,samples=180] {cos(deg((0.4*pi+2*pi)*x))};
\addlegendentry{fast: $\Omega=2.4\pi$}
% shared samples
\addplot[specred,only marks,mark=*,mark size=1.8pt,samples at={0,1,...,8}] {cos(deg(0.4*pi*x))};
\addlegendentry{shared samples}
\end{axis}
\end{tikzpicture}
\caption{Aliasing. The slow tone ($\Omega = 0.4\pi$) and the fast tone ($\Omega = 2.4\pi$) differ
by exactly $2\pi$ per sample, so they agree at every integer tick (red dots). From the samples
alone the two are indistinguishable; the faster one is aliased onto the slower.}
\label{ch04:fig-alias}
\end{figure}

\section{A worked round trip in code}
\label{ch04:code}

Let us verify the phasor identity numerically.
The snippet adds two same-frequency sinusoids directly and via their phasors, and confirms the two
agree.

\begin{lstlisting}[language=Python,caption={Adding sinusoids: time-domain vs.\ phasor.},captionpos=b]
import numpy as np

w  = 2*np.pi*5.0          # 5 Hz, angular
t  = np.linspace(0, 1, 2000)

# two sinusoids of the SAME frequency
A1, p1 = 2.0,  0.0
A2, p2 = 1.4,  np.pi/2

x_direct = A1*np.cos(w*t + p1) + A2*np.cos(w*t + p2)

# phasor sum: add complex amplitudes, then read magnitude and angle
X = A1*np.exp(1j*p1) + A2*np.exp(1j*p2)
A3, p3 = np.abs(X), np.angle(X)
x_phasor = A3*np.cos(w*t + p3)

print(f"A3 = {A3:.4f}, phi3 = {p3:.4f} rad")     # -> 2.4413, 0.6107
print("max error:", np.max(np.abs(x_direct - x_phasor)))  # ~1e-15
\end{lstlisting}

The maximum difference is at the level of floating-point round-off: the two computations are the
same operation seen from two sides. The phasor route did the hard part -- finding $A_3$ and
$\phi_3$ -- with a single complex addition.

\begin{labbox}
Modify the snippet to add \emph{three} sinusoids and reproduce Example~\ref{ch04:ex-three-sum}
(expect $A_3 = 1$, $\phi_3 = 0$). Then deliberately break the assumption: give the second sinusoid
a \emph{different} frequency ($6$ Hz instead of $5$). Plot the sum. You will see it is no longer a
clean sinusoid -- it beats and wanders -- because the phasor trick requires a common $\omega$. That
failure motivates Fourier analysis: when frequencies differ, you cannot collapse
to one arrow; you must keep them all.
\end{labbox}

\section{Toward the wrist tracker}
\label{ch04:goal}

\begin{goalbox}
Why does a chapter on sinusoids matter for a wrist-worn gesture controller that will output MIDI?
Because nearly every signal the device handles is best understood as a superposition of rotating
phasors.

\textbf{Tremor and motion.} When the forearm oscillates -- a deliberate shake, a vibrato gesture,
or involuntary physiological tremor near $8$--$12$ Hz -- the accelerometer and gyroscope axes each
read an approximate sinusoid. The \emph{same} physical oscillation appears on the $x$, $y$, and $z$
axes with different amplitudes and \emph{phases}, exactly the same-frequency-different-phase
situation of Section~\ref{ch04:phasors}. Combining the axes is phasor addition; the phase
differences between axes encode the \emph{direction} of the motion in space.

\textbf{The frequency response is your filter.} The eigenfunction property of
Section~\ref{ch04:eigen} is the reason we can design a digital filter to keep the slow,
intentional gesture band and reject the fast tremor band: each frequency component is multiplied by
its own gain $\abs{H(\omega)}$, independently. We will lean on this constantly.

\textbf{A MIDI vibrato.} If you map the \emph{frequency} of a detected wrist oscillation to a MIDI
modulation-wheel value and its \emph{amplitude} to depth, a performer can ``dial in'' vibrato by
shaking faster or wider. Estimating that single frequency and amplitude amounts geometrically to finding the
dominant phasor in the signal -- the simplest possible Fourier analysis.
\end{goalbox}

\begin{goalbox}
\textbf{A real limit, stated plainly.} A wrist- or forearm-worn IMU measures the motion of the
\emph{whole limb}. It can sense gross orientation, swings, taps, and oscillations like the tremor or
vibrato above. It \emph{cannot}, by itself, recover the continuous joint angles of individual
fingers: bending only the index finger produces almost no signal at the wrist. Finger-level
intent requires a different physics -- surface EMG or force myography (FMG) reading the muscle
activity in the forearm (Chapter~37), and continuous finger \emph{pose} requires sensors on the
fingers themselves, such as instrumented rings. Treat the IMU as a detector of coarse, often
discrete, whole-hand gestures, and reach for EMG/FMG when you need to know what the fingers are
doing. We will keep returning to this boundary so that nothing we build over-promises. This
framing is well supported by the wearable-sensing literature (surface-EMG gesture studies, and the
known limitations of inertial-only hand tracking).
\end{goalbox}

\begin{buildbox}
A concrete first experiment. Strap an MPU-6050 IMU to your forearm, sample the accelerometer at,
say, $200$ Hz, and shake your hand back and forth at a steady rate. Then estimate the oscillation
frequency from the zero-crossings of one axis -- the rate at which the cosine's shadow crosses its
center line is twice the frequency, since each cycle has two crossings.

\begin{lstlisting}[language=C,caption={Crude zero-crossing frequency estimate on a microcontroller.},captionpos=b]
// ax[]: N recent accelerometer samples, mean-removed; fs = sample rate (Hz)
int crossings = 0;
for (int n = 1; n < N; n++) {
    if ((ax[n-1] <= 0.0f && ax[n] > 0.0f) ||   // upward crossing
        (ax[n-1] >= 0.0f && ax[n] < 0.0f))     // downward crossing
        crossings++;
}
float freq_hz = (0.5f * crossings) * (fs / (float)N);  // cycles per second
\end{lstlisting}

This is the cheapest possible frequency estimator and it works for a clean single tone. It is
fragile against noise and multiple tones, which is precisely why we will build the Fourier
transform: to read \emph{all} the frequencies at once, robustly. But feeling a number come out that
tracks your shaking is a good first taste of turning motion into a control signal.
\end{buildbox}

\section{History and a note on notation}
\label{ch04:history}

\begin{historybox}
The sinusoidal foundation of physics and engineering took shape through a fifty-year debate
over music and mathematics. In 1746, \textbf{Jean le Rond d'Alembert} derived the one-dimensional
wave equation $\partial^2 y/\partial t^2 = c^2 \partial^2 y/\partial x^2$ for a vibrating string and
solved it using traveling waves $y(x,t) = f(x-ct) + g(x+ct)$. In 1748, \textbf{Leonhard Euler}
published his famous identity $e^{\jj\theta} = \cos\theta + \jj\sin\theta$ in his \emph{Introductio
in analysin infinitorum} and derived the standing-wave sinusoidal modes of the string.

Then, in 1753, \textbf{Daniel Bernoulli} asserted that \emph{every}
possible motion of a musical string, no matter how it was initially plucked, could be expressed
as a simultaneous superposition of sinusoidal harmonics:
$y(x,t) = \sum_{n=1}^\infty A_n \sin\big(\frac{n\pi x}{L}\big)\cos\big(\frac{n\pi c t}{L}\big)$.
Both d'Alembert and Euler objected. A plucked string starts with a sharp triangular kink
(a point with no derivative). Euler argued passionately that because every sine wave is smooth and
infinitely differentiable, no infinite sum of smooth sinusoids could ever represent a kinked line or
an ``arbitrary curve drawn by a free stroke of the hand.'' Bernoulli possessed the correct physical
acoustical intuition, but lacked the tool of orthogonal integration to compute the coefficients
$A_n$. The controversy remained unresolved until Fourier's work in 1807.

The practical use of sinusoids in electrical engineering arrived in 1893 with \textbf{Charles
Proteus Steinmetz}. A brilliant, hunchbacked German immigrant and socialist exile working at General
Electric, Steinmetz was confronted with the mathematical difficulty of designing alternating-current
(AC) power networks. Analyzing circuits with resistors, inductors, and capacitors required solving
intricate systems of coupled differential equations with tedious trigonometric identities. In an
influential 1893 paper presented to the American Institute of Electrical Engineers, Steinmetz
showed that because $e^{\jj\omega t}$ is an eigenfunction of differentiation, one could freeze the
common time factor $e^{\jj\omega t}$ and represent steady-state voltages and currents as static
complex numbers: \emph{phasors} $V = V_0 e^{\jj\phi}$. By assigning inductors an impedance
$Z_L = \jj\omega L$ and capacitors $Z_C = 1/(\jj\omega C)$, differential equations collapsed into
elementary high-school Ohm's law algebra: $V = I Z$. Steinmetz's algebraic phasor method made
complex numbers a standard tool in electrical engineering and supported the design of the Niagara Falls
hydroelectric plant and the continental power grid.
\end{historybox}

\begin{marginnotebox}
A notation caution for newcomers crossing between fields. Mathematicians and physicists write the
imaginary unit as $i$ and often the complex exponential as $e^{-i\omega t}$ (note the minus sign);
electrical engineers write $\jj$ and frequently $e^{+\jj\omega t}$. The sign convention only flips
the sense of rotation and the sign of phase; nothing physical changes, but you must pick one
convention and hold it. This book uses $\jj$ and the $e^{+\jj\omega t}$ convention throughout.
\end{marginnotebox}

\begin{aestheticbox}
One geometric change simplified a great deal. We started with three real numbers ($A$, $\omega$,
$\phi$) and a stubborn trig-identity problem. By recognizing the cosine as the shadow of a
rotation, we compressed amplitude-and-phase into a \emph{single} complex number, turned adding
waves into adding arrows, and -- almost for free -- discovered that these same complex
exponentials are the eigenfunctions that make every linear system simple. The single idea of
rotation connects the algebra, the signals, and the systems, and the rest of the book relies on
that connection.
\end{aestheticbox}

\section*{Summary}
\addcontentsline{toc}{section}{Summary}

\begin{itemize}
  \item A real sinusoid $A\cos(\omega t + \phi)$ is fixed by three numbers: amplitude (vertical
        scale), angular frequency (how fast, $\omega = 2\pi f$), and phase (sideways shift /
        starting angle).
  \item A cosine is the horizontal shadow of a point rotating at constant speed; a sine is the
        vertical shadow. The wiggle's apparent acceleration is an artifact of projection.
  \item Euler's formula $e^{\jj\theta} = \cos\theta + \jj\sin\theta$ makes a rotating vector a
        single number. The complex exponential $A e^{\jj\phi} e^{\jj\omega t}$ packages all three
        parameters; the constant $X = Ae^{\jj\phi}$ is the \emph{phasor}.
  \item Same-frequency sinusoids add by adding phasors tip-to-tail, replacing trig identities with
        vector geometry.
  \item A real cosine is two counter-rotating half-arrows at $\pm\omega$; this is what ``negative
        frequency'' means and why real spectra are symmetric.
  \item Complex exponentials are the eigenfunctions of every LTI system: in, $H(\omega)$ times the
        same exponential, out. This is the reason we decompose signals into sinusoids at all.
  \item Discrete-time exponentials hop around the circle; frequency becomes periodic with period
        $2\pi$, foreshadowing aliasing and Nyquist.
\end{itemize}

\begin{connectionbox}
\textbf{Where We Are and Where We Go Next}
\begin{itemize}
  \item \textbf{What we have so far}: Viewing sinusoids as shadows of rotating arrows $e^{\jj\omega t}$
  turns trig identities and derivatives into complex multiplication. Adding same-frequency tones
  becomes tip-to-tail vector addition, time derivatives become multiplication by $\jj\omega$, and
  linear time-invariant systems scale each $e^{\jj\omega t}$ by one complex number $H(\omega)$.
  \item \textbf{What is still missing}: A single phasor describes only one steady frequency running
  forever from $t = -\infty$ to $+\infty$. By itself it cannot describe a chord, a drum strike, a
  spoken vowel, or a quick wrist flick -- and two phasors at \emph{different} frequencies cannot be
  frozen into one static arrow because one keeps spinning relative to the other.
  \item \textbf{Next}: Chapter~5 puts multiple frequencies together and introduces the Fourier
  winding machine to separate them. Chapter~6 builds the Fourier series for periodic waveforms,
  Chapter~7 stretches the period to infinity to handle one-shot transients, and Chapter~10 slides a
  short window through time to track frequencies as they change.
\end{itemize}
\end{connectionbox}

\section{Exercises}
\label{ch04:exercises}

\begin{exercise}[Units and conversions]
A signal completes $440$ cycles per second (the musical note A4). Give its frequency $f$, angular
frequency $\omega$, and period $T$. If you sample it at $f_s = 8000$ Hz, what is the discrete
angular frequency $\Omega = \omega/f_s$ in radians per sample? Is it below the Nyquist limit
$\pi$?
\end{exercise}

\begin{exercise}[Cosine, sine, and phase]
Express $\sin(\omega t)$, $-\cos(\omega t)$, and $-\sin(\omega t)$ each in the form
$\cos(\omega t + \phi)$ with $\phi \in (-\pi, \pi]$. Sketch the four phasors $1$, and the three you
found, as arrows in the complex plane and confirm they sit at $0$, $-\pi/2$, $\pi$, $\pi/2$.
\end{exercise}

\begin{exercise}[Phasor addition by hand]
Add $3\cos(\omega t) + 4\cos(\omega t - \tfrac{\pi}{2})$ using phasors. Find the resultant
amplitude and phase, and write the answer as a single cosine. (Hint: the phasors are $3$ and
$-4\jj$; the geometry is a 3--4--5 right triangle.)
\end{exercise}

\begin{exercise}[When cancellation is total]
For what phase difference $\Delta\phi = \phi_2 - \phi_1$ do two \emph{equal-amplitude},
same-frequency sinusoids cancel exactly to zero? Explain with the arrow picture, and give the
amplitude of the sum as a function of $\Delta\phi$ in general. (You should find the sum amplitude
is $2A\abs{\cos(\Delta\phi/2)}$.)
\end{exercise}

\begin{exercise}[Euler both ways]
Using $\cos\theta = (e^{\jj\theta} + e^{-\jj\theta})/2$, derive the identity
$\cos^2\theta = \tfrac12 + \tfrac12\cos(2\theta)$ by multiplying out the exponentials. Notice you
never needed to remember a half-angle formula -- the exponentials did it for you.
\end{exercise}

\begin{exercise}[Product to sum]
Show that $\cos(\omega_1 t)\cos(\omega_2 t) = \tfrac12\cos((\omega_1-\omega_2)t) +
\tfrac12\cos((\omega_1+\omega_2)t)$ by writing each cosine with Euler's formula and expanding. This
``beating'' of two close frequencies is the basis of amplitude modulation (Chapter~23) and of the
amplitude wobble you saw in the lab exercise.
\end{exercise}

\begin{exercise}[The eigenfunction argument]
Suppose an LTI system has output $y(t) = \mathcal{H}\{e^{\jj\omega t}\}$. Using only linearity and
time-invariance, fill in each step of the derivation in Section~\ref{ch04:eigen} that
$y(t) = H(\omega)e^{\jj\omega t}$. Where exactly is time-invariance used, and where is linearity?
\end{exercise}

\begin{exercise}[Frequency response of a delay]
A pure delay system outputs $y(t) = x(t - \tau)$. Show it is LTI, and by feeding in
$x(t) = e^{\jj\omega t}$ find its eigenvalue $H(\omega)$. Confirm $\abs{H(\omega)} = 1$ (a delay
does not change amplitude) and $\arg H(\omega) = -\omega\tau$ (it adds a phase proportional to
frequency -- ``linear phase'').
\end{exercise}

\begin{exercise}[Discrete periodicity]
Show that the discrete sequences $x_1[n] = \cos(0.3\pi\, n)$ and $x_2[n] = \cos(2.3\pi\, n)$ are
identical for all integers $n$. Which one corresponds to a frequency inside the fundamental band
$(-\pi, \pi]$? What does this say about a continuous tone that, when sampled, produces these
numbers?
\end{exercise}

\begin{exercise}[Three-axis tremor, conceptual]
A forearm tremor at $f = 9$ Hz appears on the three accelerometer axes as
$a_x = 0.20\cos(\omega t)$, $a_y = 0.15\cos(\omega t - 1.2)$, $a_z = 0.05\cos(\omega t + 0.4)$ (all
in $g$, with $\omega = 2\pi\cdot 9$). Using phasors, find the amplitude and phase of the combined
projection $a_x + a_y + a_z$. Then explain in one sentence why the \emph{phase differences} between
axes carry information about the direction of the shaking that the amplitudes alone do not.
\end{exercise}

\begin{exercise}[Build reflection]
The zero-crossing estimator in the build box reports half the crossing count times $f_s/N$. Explain
why a single noisy sample near a zero crossing can cause it to badly over-count, and propose one
simple fix (e.g.\ a small hysteresis band or a low-pass pre-filter). Connect your answer to why we
will prefer Fourier-based frequency estimation later in the book.
\end{exercise}
